\documentclass[a4paper]{article}
\usepackage[pdftex]{graphicx}
\usepackage[utf8]{inputenc}
\usepackage{enumerate}
\usepackage{amssymb}
\usepackage{href-ul}
\hypersetup{
	colorlinks=true,
	linkcolor=blue,
	urlcolor=blue}
\usepackage{geometry}
\geometry{a4paper, top=15mm, left=15mm, right=15mm, bottom=15mm,
	headsep=10mm, footskip=12mm}

\begin{document}
{\bf Magnetismus}
\begin{enumerate}[1.]

%Aufgabe 1
{\bf \item Stromdurchflossene Leiter im Magnetfeld}\\
Fülle die Lücken richtig aus:\\
\renewcommand{\baselinestretch}{2}\normalsize
Auf einen stromdurchflossenen Leiter im Magnetfeld wirkt die so genannte \rule{70 pt}{0.5 pt}.
Ihre Richtung kann verändert werden durch eine Richtungsänderung des \rule{70 pt}{0.5 pt}oder des \rule{70 pt}{0.5 pt}. Man beobachtet keine Wirkung auf den Stromleiter im Magnetfeld, wenn Leiter und Magnetfeldlinien \rule{140 pt}{0.5 pt}. Der Ausschlag des stromdurchflossenen Leiters im Magnetfeld ist umso größer, je:


\begin{itemize}
\item \rule{140 pt}{0.5 pt}
\item \rule{140 pt}{0.5 pt}
\item \rule{140 pt}{0.5 pt}
\end{itemize}

\renewcommand{\baselinestretch}{1}\normalsize

%Aufgabe 2
{\bf \item Aluminiumstange auf Eisenschienen}\\
\begin{minipage}{0.45\textwidth}
Eine Aluminiumstange wird auf zwei Metallschienen gelegt, sodass sie frei rollen kann. Die Anordnung befindet sich zwischen den Polschuhen eines Hufeisenmagneten wie in der Skizze dargestellt.
\end{minipage}
\hfill
\begin{minipage}{0.5\textwidth}
\includegraphics[width=7 cm]{magnschi132}
\end{minipage}

\begin{enumerate}[a)]
\item Erläutere kurz was man beobachtet, wenn der Stromkreis geschlossen wird (Begründung).
\item Was bewirkt eine Umpolung der Stromrichtung ?
\end{enumerate}

%Aufgabe 3
{\bf \item Elektromotor}

\includegraphics[width=13 cm]{mgnlsch132}

\begin{enumerate}[a)]
\item Zeichne die Kraftpfeile in Stuation I und II ein und kennzeichne die Drehrichtung der Leiterschleife.
\item Trage die Polung der Leiterschleife in Situation III ein, so dass die Drehrichtung erhalten bleibt.
\item Erläutere die Funktionsweise des Elektromotors, verwende und erkläre dabei die Begriffe Rotor, Stator, Totpunkt und Polwender.   
\end{enumerate}

\begin{minipage}{0.55\textwidth}
{\bf \item Eine Leiterschleife soll sich in die angegebene Richtung bewegen.} Wie muss die Stromquelle gepolt sein? Erläutere deine Überlegungen unter Anwendung der Drei-Finger-Regel.
\end{minipage}
\hfill
\begin{minipage}{0.4\textwidth}
\includegraphics[width=5 cm]{leischau214}
\end{minipage}

{\bf \item Was versteht man unter der Lorentzkraft?}

{\bf \item Die von einem Magnetfeld hervorgerufene Kraftwirkung auf freie Ladungsträger} kann zum Beispiel mit Hilfe eines Fadenstrahlrohres gezeigt werden. Beschreibe auch anhand einer beschrifteten Skizze die Funktionsweise eines Fadenstrahl\-rohrs genau.

\begin{minipage}{0.3\textwidth}
{\bf \item Ein Elektron tritt in ein räumlich begrenztes homogenes Magnetfeld ein.} Zeichne den wei\-teren Bahnverlauf sowie exemplarisch einige auf das Elektron wirkende Kräfte ein. Erkläre deine Lösung.
\end{minipage}
\hfill
\begin{minipage}{0.65\textwidth}
\hspace{30 pt}
\includegraphics[width=7 cm]{elo214}
\end{minipage}


\end{enumerate} 
\fbox{
	\begin{minipage}{0.5\textwidth}
	Weitere Arbeitsblätter gibt es unter 
	
	\href{https://www.okuyakl.de}{www.okuyakl.de}
	\end{minipage}
	\hfill
	\begin{minipage}{0.4\textwidth}
		\includegraphics[width=1.5 cm]{../../viecher/zwe03}
		\includegraphics[width=2 cm]{../../viecher/afanticon1}
		
\end{minipage}}

\end{document}

